	\subsubsection{Beliefs over Returns to College Degree}\label{sec:additional_belief_data}
		Our survey included the survey instruments developed in \citealp{attanasio2014education} to elicit students beliefs about returns to a college degree. We elicited beliefs about the distribution of wages at age 30 with and without a college degree. We find that students think that, on average, the return to a college degree is 82 percent. This is in line with observed differences in wages between Chileans with and without a college degrees, and in line with results from other surveys on different samples of Chilean high-school students (\citealp*{hastings2016informed}).\par 
		
		We found that the treatment did not have any impact on student beliefs about returns to education (no impacts on the mean nor on the variance of the returns), as reported in Table \ref{expearnings}.
		
		\input{output/tables/exp_earnings}

		Expected (mean) earnings were directly elicited, and we also estimated them, together with the variance of earnings, from elicited c.d.f. values \citep{miglino_tincani_subjective_earnings_2026}. We report results on both measures of expected earnings, for comparison.\par
		The survey questions asked “How much do you expect to earn per month with (without) a college degree on average?” and “How likely are you to earn at least X pesos per month with (without) a college degree?” where X={$200,000$, $800,000$} without a degree and  X={$300,000$, $1,000,000$} with a degree. To calculate the mean and variance of expected earnings using the answers to these questions, we fit the reported c.d.f. values using log-normal distributions for each respondent in the sample. In the estimation sample we kept only the students that answered at least two questions for each scenario (with and without a degree), because we needed at least two c.d.f. values to estimate the mean and variance of the Log-normal distribution. Finally, we used the Generalized Method of Moments to find the mean and variance of the log-normal distribution that minimize the distance of the simulated mean and simulated c.d.f. values from their data analogues.\par
		For variance regressions we use median regressions because the variance is very vulnerable to outlier survey responses in which a student gives the same probability to the likelihood that his/her earnings at age $30$ will be above two different values. 
	
